\section{DeltaTrace}
\label{sec:method}

\subsection{The counterfactual question and attribution target}
We explain how masking selected prompt tokens changes the score of a fixed target span. The inputs are a trained language model with fixed parameters $\theta$, a prompt $x$ containing the context and question, and a stored response $y=(y_1,\ldots,y_L)$ of $L$ tokens. Let $\mathcal{T}$ be the positions of the target span within $y$. Both executions score these positions using the same preceding response tokens. The scalar target score is
\begin{equation}
F(x;y)=\sum_{t\in\mathcal{T}}\log p_\theta(y_t\mid x,y_{<t}),
\label{eq:target}
\end{equation}
where $p_\theta$ is the model's next-token probability and $y_{<t}$ is the response prefix before position $t$. A higher score means the model assigns greater probability to the target span given its preceding tokens. For full-response attribution, $y$ includes the reasoning, answer, and end-of-sequence (EOS) token, and $\mathcal{T}=\{1,\ldots,L\}$. The response tokens and target positions stay fixed, while their hidden representations can change as they read the different prompts.

The source positions are the context or question tokens chosen for explanation. Let $x_1$ be the original prompt. Its reference $x_0$ masks all source tokens together with EOS at the same positions (Figure~\ref{fig:overview}(a)), keeping all other prompt tokens fixed. Both executions process the full input sequence and score the same target span, giving
\[
\D F=F(x_1;y)-F(x_0;y).
\]
A positive $\D F$ means masking these tokens lowers the target score. The attribution task is to divide this difference among the source tokens. Each token receives a signed contribution $A_i$, and adding the contributions recovers $\D F$. Positive and negative scores indicate support and opposition under this reference (Figure~\ref{fig:overview}(d)).

Subscripts $0$ and $1$ denote reference and original activations, and $\D U=U_1-U_0$ denotes their difference. For source token $i$, $e_{i,0}$ is the EOS embedding and $e_{i,1}$ is its original embedding, so $\D e_i=e_{i,1}-e_{i,0}$. Tokens held fixed have zero embedding difference. The reverse pass converts these embedding differences into token contributions.

\subsection{From the score difference to token contributions}
\label{sec:reverse}
\dt{} replaces tangent slopes with secant slopes to trace the measured output change back to the inputs. For a scalar activation, the secant connects its reference and original points: its slope converts the horizontal input displacement into the vertical output displacement. A tangent slope gives only a first-order approximation to this finite change. Equal inputs use the local derivative.

The reverse pass carries $m$ for each activation entry, its effective secant slope to the final response score under the chosen allocation rules. Multiplying this slope by the entry's displacement gives its assigned contribution. Subscripts identify the associated activation. The pass starts with $m_F=1$, since the full score difference is to be allocated. Table~\ref{tab:reverse-comparison} contrasts tangent and secant propagation through the same activation.

\begin{table}[htbp]
\centering
\small
\caption{\textbf{The same activation, two backward calculations.} The activation has input $a$ and output $h=f(a)$. Here $d$ denotes an infinitesimal change and $\Delta$ the observed difference between the two executions.}
\label{tab:reverse-comparison}
\renewcommand{\arraystretch}{1.35}
\begin{tabular}{@{}lcc@{}}
\toprule
 & \shortstack{Ordinary backpropagation\\Tangent} & \shortstack{\dt{}\\Secant} \\
\midrule
Local slope & $f'(a_1)$ & $\dfrac{\D h}{\D a}=\dfrac{h_1-h_0}{a_1-a_0}$ \\
\addlinespace[3pt]
Output change & $dh=f'(a_1)\,da$ & $\D h=\dfrac{\D h}{\D a}\,\D a$ \\
\addlinespace[3pt]
Backward calculation & $\dfrac{\partial F}{\partial a}=f'(a_1)\dfrac{\partial F}{\partial h}$ & $m_a=\dfrac{\D h}{\D a}\,m_h$ \\
\bottomrule
\end{tabular}
\end{table}

For array operations, we use the operation's formula and both executions to express output changes as weighted sums of input changes. These local slopes form the matrix $D_f$. Linear operations use their existing weights. For a product, splitting the term involving both input changes equally between them gives equal weight to the two possible change orders. This produces the mean attention weights and mean values in Section~\ref{sec:attention}.

The same local map relates observed differences and propagates $m$ backward through its transpose,
\begin{equation}
\D h=D_f\D a,
\qquad m_a=D_f^\top m_h.
\label{eq:reverse-contribution}
\end{equation}
The transpose multiplies each output's effective slope by the corresponding local slopes and sums the products for each input. This expresses the same contribution in terms of input changes. The local slopes replace derivatives, and the propagated slopes replace gradients. The two products have the forms of a Jacobian--vector product (JVP) and a vector--Jacobian product (VJP), respectively.\footnote{PyTorch documentation for \href{https://docs.pytorch.org/docs/2.14/generated/torch.func.jvp.html}{\texttt{torch.func.jvp}} and \href{https://docs.pytorch.org/docs/2.14/generated/torch.func.vjp.html}{\texttt{torch.func.vjp}}.}

At the input embeddings, the propagated $m_{e_i}$ give the effective secant slopes for each token. Summing their products with the corresponding embedding displacements yields the token contributions,
\begin{equation}
A_i=\ip{m_{e_i}}{\D e_i},
\qquad \sum_i A_i=\D F.
\label{eq:attribution}
\end{equation}
Every backward calculation preserves the contribution being divided, so the token scores add to the measured response-score difference. Positive and negative scores identify support and opposition across the chosen response span. They use natural-log units (nats) and can be summed into name, sentence, or passage contributions. A fixed token's embedding does not change, so its contribution is zero. Appendix~\ref{app:conservation} gives the full derivation.

We next apply these rules to attention and Gated DeltaNet (Figure~\ref{fig:overview}(b, c)). Appendix~\ref{app:operators} gives the formulas for linear layers, additions, activation functions, and normalizations in the same reverse pass.

\subsection{Attention example}
\label{sec:attention}
For attention, $P_{ij}$ is the weight with which position $i$ reads source $j$, and the rows of $V$ contain source value vectors. The output $Y=PV$ is their weighted sum, with $P$ obtained by softmax normalization of query--key similarity scores.

Let $\D P$, $\D V$, and $\D Y$ denote the observed differences in weights, values, and output between the two executions. Expanding the product gives a term involving both input changes. DT splits this term equally between the inputs, giving equal weight to the two possible change orders,
\begin{equation}
\begin{aligned}
\D Y
&=P_0\D V+\D P\,V_0+\underbrace{\D P\,\D V}_{\text{split equally}}\\
&=\underbrace{\left(P_0+\tfrac12\D P\right)}_{\bar P}\D V
+\D P\underbrace{\left(V_0+\tfrac12\D V\right)}_{\bar V}.
\end{aligned}
\label{eq:pv}
\end{equation}
The terms in parentheses are the means of the original and reference weights and values. These are the two sets of local slopes represented by $D_f$ in Equation~\ref{eq:reverse-contribution}. The mean weights multiply the value changes, and the mean values multiply the weight changes.

The incoming slopes $m_Y$ weight the two terms in the second line of Equation~\ref{eq:pv}, dividing the output's assigned contribution between value and weight changes. The backward calculation uses the same two products as ordinary backpropagation, with means replacing the original activations,
\begin{equation}
\begin{aligned}
\text{Backpropagation}\quad &
\frac{\partial F}{\partial V}=P_1^\top\frac{\partial F}{\partial Y},
&\quad \frac{\partial F}{\partial P}=\frac{\partial F}{\partial Y}V_1^\top,\\
\text{DeltaTrace}\quad &
m_V=\bar P^\top m_Y,
&\quad m_P=m_Y\bar V^\top.
\end{aligned}
\label{eq:pv-reverse}
\end{equation}
The returned slopes continue through preceding operations to the input embeddings.

Softmax makes each attention row sum to one, so increasing one source's weight reduces the total assigned to the others. The backward calculation for softmax carries this competition through a weighted row average subtracted from the incoming slopes. Its weights depend on the probabilities in both executions. The slopes then pass through the query--key scores to their input representations. The same construction handles the vocabulary normalization used to score each response token. Appendix~\ref{app:softmax} gives the formula and its derivation.

\subsection{Gated DeltaNet example}
\label{sec:memory}
Gated DeltaNet implements linear attention through recurrent memory that stores key--value associations across token positions \citep{yang2024gated}. Figure~\ref{fig:overview}(c) shows one step from the previous memory $S_{t-1}$ to the updated memory $S_t$ and its readout $o_t$. A retention gate controls how much earlier memory is kept, and a write gate controls a correction toward the current value at its key. A query reads from $S_t$ to produce $o_t$. Normalization, an output gate, and a linear projection then produce the layer output. Appendix~\ref{app:memory} gives the state equations.

Information written by an earlier token can affect a later readout, so its contribution must pass backward through the intervening memory updates. The readout slopes $m_{o_t}$ pass backward to the memory and query. At $S_t$, the contribution from this readout joins those from later memory states. The resulting $m_{S_t}$ passes through the update to $S_{t-1}$, the value, the key, and the retention and write gates (Figure~\ref{fig:overview}(c)). The returned slopes continue through preceding operations to the input embeddings, where Equation~\ref{eq:attribution} gives token scores.

Within each GDN layer, the original and reference activations define two decompositions of the same readout changes over the complete memory sequence. Both backward calculations receive the same readout slopes. Their returned slopes for each value, key, query, and gate are averaged with equal weight before propagation continues to earlier layers. This averaging stays within the model's single reverse pass. The separate output-gating product also splits contributions between its two changing inputs using their two-run means. Appendix~\ref{app:memory} gives these calculations.
